\section{Prospective Execution-Path Prediction}
\label{sec:results}
\label{sec:diagnoses}
% B570 prediction summary from frozen discovery and reserved validation.
% This replaces the generic evaluation-flow graphic; no new result is inferred.
\begin{figure*}[t]
\centering
\small
\begin{tabular}{@{}p{0.95in}p{2.12in}p{3.08in}@{}}
\toprule
Stage & Design or comparator & Complete generated-ID endpoint \\
\midrule
Discovery & 12 prompts $\times$ 3 runs; select target-only UB1 as H1. & UB1 predicts the serialized trajectory in \textbf{36/36} runs. \\
Hypothesis freeze & After discovery; fix H1, reference, tasks and endpoint. & Before reserved reference/predictor inference; H2 remains unrun. \\
Prediction freeze & After pre-live inference; seal all 72 UB1 arrays. & Before the first reserved live outcome; no prediction edits. \\
Live test & Compare each of 72 live serialized runs with both frozen paths. & UB1 predictor/live \textbf{72/72}; UB128 reference/live \textbf{60/72 (BLOCKED)}. Four prompts (12 runs) are reference-divergent. \\
Divergent case & HumanEval/134; first difference at generated-ID index 88. & Reference ID 914 $\rightarrow$ predictor/live ID 1946 in all 3 runs. \\
\bottomrule
\end{tabular}
\caption{\textbf{B570 prospective prediction.} Hypothesis and full-array freezes are separate stages. The 24 prompts $\times$ three repeat-stable runs test complete capped trajectories, including four non-reference predictions. Reported ordering and retained arrays are reviewed in the appendix; prediction success does not restore qualification.}
\label{fig:translation}
\end{figure*}

\subsection{B570: a draft-free full-trajectory predictor}
\label{sec:b570mechanism}
A target-only harness first reproduces 90/90 UB128 reference sequences. On 12 exposed discovery tasks, H1 target-only UB1 matches 36/36 native serialized trajectories. The frozen decision tree selects H1 and leaves H2 unrun. Conservative exposure exclusion then selects 24 reserved tasks. H1 freezes before reserved reference/predictor inference, and all 72 complete predictions freeze before reserved native live (Figure~\ref{fig:translation}).

Prediction/live agrees 72/72 over three stable runs per prompt, including all 12 runs of four reference-divergent prompts. The same-set UB128 reference agrees only 60/72 and remains \emph{BLOCKED}. Every reserved array reaches 128 generated IDs. Any wrong later ID would falsify the prediction, rather than merely reproduce a first-difference signature. The four first differences occur on /134 at index88 (914$\to$1946), /42 at 73 (432$\to$1817), /145 at 103 (1184$\to$686), and /138 at 93 (14374$\to$9064); all later IDs are also predicted. Appendix J retains every committed array and freeze role.

This answers a question that qualification alone cannot. Reference mismatch establishes that the path is unsuitable for strict replacement; it does not establish whether a target-only program can describe its full output. Here the selected target-only UB1 program suffices without receiving draft weights, proposals or future live IDs. That result rules out the need for those inputs \emph{in this successful predictive description}. It does not show that drafting is causally irrelevant inside native execution, that UB1 is the unique predictor, or that the output contract is restored.

\subsection{Sampled numerical sensitivity is not an immediate split}
A separate frozen fork study uses 24 new prompts $\times$ three, sampling the first verification block after at least 64 committed IDs. UB1/UB128 branches share checked KV, prefix, positions, proposals and cloned sampler state. Scores change at \FUForkScores/72 forks, but raw top-one and actual immediate decisions change at \FUForkImmediate/72. The immediate comparison contains 195 reachable paired decisions; 216 raw-logit rows are correlated within prompts and blocks.

\begin{samepage}
Existing-record descriptors put these results on a score scale. Median absolute top-one score change is \DerScoreAbsMedian, with p95 0.3437 and maximum 2.3503. Raw top-two margins have minima 0.2291 (UB1) and 0.2003 (UB128), and medians 17.5021 and 17.5128. These are marginal descriptors, not a certified decision-boundary distance or a uniform perturbation bound. In particular, the study does not supply a sampled near-tie witness. Changed scores with stable sampled decisions show local numerical sensitivity and immediate stability, not a whole-generation null effect.
\end{samepage}

The alternate branch restores draft state and continues with UB128 target blocks; later proposals follow its altered live history. Complete trajectories differ in \FUForkLater/72 runs on \FUForkLaterPrompts/24 prompts, first diverging 10--50 positions after the fork. Persistent continuation therefore confounds attributing a delayed difference to one sampled block. The study neither explains the original H1 four tasks nor establishes immediate reference restoration. Appendices L/P retain score/margin distributions, all six later-sensitive prompts and clone/R0 controls.

\subsection{RTX: discriminating predefined candidates on one set}
\label{sec:rtxreplay}
Earlier RTX H1 fails exposed discovery at 24/36 and stops with zero reserved runs. A separate V2 study obtains 72/72 matches from proposal-conditioned H2 and fresh R\_ONLY but lacks a same-set target-only competitor (Appendix M). V4 addresses that missing comparison; it is not a retest of one portable B570 predictor.

All 19 remaining exposure-audited tasks enter V4. Four pre-live conditions---UB1 reference, target-only UB2, H2 and fresh R\_ONLY---produce 228 complete arrays. All arrays/tapes are sealed and remotely retrieved before 57 native batched live runs. H2's draft worker P uses its own earlier target feedback to construct proposals; target worker R computes its own decisions. Fresh R\_ONLY reads the sealed tape sequentially, and native live imports no tape. This separates application/API orchestration on shared llama.cpp/GGML/CUDA numerical routines (Table~\ref{tab:inputcontracts}), not adversarial OS access or an independent numerical implementation.

Table~\ref{tab:rtxv4trajectory} shows complete-array agreement and the cross comparison hidden by aggregate totals. Target-only UB2/live is 51/57 on 17/19 prompts. H2/live, R\_ONLY/live and H2/R\_ONLY are each 57/57; both programs predict the two reference-divergent trajectories. All five conditions are repeat-stable. UB2 and reference both score 51/57, but /18 refutes only the reference and /25 only UB2; /15 refutes both.
\begin{table}[t]
\caption{RTX V4 complete prediction/live agreement on 19 prompts $\times$ three. Separate P/R orchestration uses shared numerical routines. Below: full-array failures on all distinguishing tasks, stable in three repeats; replay matches every task. These are prediction comparisons, not qualification.}
\label{tab:rtxv4trajectory}
\centering
\begin{tabular}{@{}lrr@{}}
\toprule
Predictor / comparator & Pairs / 57 & Prompts / 19 \\
\midrule
Target-only UB2 & 51 & 17 \\
Separate P/R (H2) & 57 & 19 \\
Fresh R\_ONLY & 57 & 19 \\
UB1 reference & 51 & 17 \\
\bottomrule
\end{tabular}
\par\smallskip
\begin{tabular}{@{}llll@{}}
\toprule
Task & UB2 & Reference & H2 / R\_ONLY \\
\midrule
/15 & Fail & Fail & Match / Match \\
/18 & Match & Fail & Match / Match \\
/25 & Fail & Match & Match / Match \\
Other 16 & Match & Match & Match / Match \\
\bottomrule
\end{tabular}
\end{table}


The prospective result rejects the specified UB2 program's full-cohort sufficiency. Matching capacity does not establish identical occupied forward shapes, inputs or prior state transitions. Replay supplies a successful richer execution description; its success does not reject every target-only hypothesis or prove proposal/history necessary, minimal or unique. Tape construction uses drafting, and replay shares numerical code. Appendix Q gives first-difference IDs, complete endpoints and the failed V3 capture history; V3 contributes no observations to these denominators.
